\begin{examplebox}
\textbf{例 1：已知一段长度，还要确定延长方向。}
在 $\triangle ABC$ 中，$AB=4,AC=6$，外角平分线交直线 $BC$ 于 $D$。
已知 $BD=8$，求 $DC$，并判断点的顺序。
\[
 \frac8{DC}=\frac46\quad\Longrightarrow\quad DC=12.
\]
因为 $AB<AC$，顺序为 $D,B,C$，所以 $BC=DC-BD=4$。
\textbf{回验：} 比例 $8/12=4/6$；三角形存在性满足 $|6-4|<4<6+4$。
故 $\boxed{DC=12}$。

\tcbbreak
\textbf{例 2：已知三边，先定位再列差。}
已知 $AB=5,AC=3,BC=4$，求外角平分线的交点到 $B$ 的距离。
由于 $AB>AC$，顺序为 $B,C,D$。设 $DC=x>0$，则 $BD=x+4$。
\[
 \frac{x+4}{x}=\frac53
 \quad\Longrightarrow\quad 3x+12=5x
 \quad\Longrightarrow\quad x=6.
\]
所以 $\boxed{BD=10}$。
\textbf{回验：} $10-6=4$，$10/6=5/3$；通式也给出 $BD=4\times5/(5-3)=10$。
这与直觉理解中的第二幅图一致。

\textbf{例 3：内外角平分线一起出现，求分点间距离。}
已知 $AB=4,AC=2,BC=3$。内角平分线交 $BC$ 于 $E$，外角平分线交直线 $BC$ 于 $D$，求 $DE$。
\begin{center}
\begin{tikzpicture}[font=\small]
\draw[black!30,rounded corners=2pt] (0,0) rectangle (14.4,8.2);
\node[anchor=west,font=\small\bfseries] at (0.5,7.72) {内分看和，外分看差：例 3 的顺序是 $B,E,C,D$};
% MODEL 4 2 3
\begin{scope}[shift={(2.4,3)},scale=1.6]
\coordinate (A) at (3.5,1.9364916731);
\coordinate (B) at (0,0);
\coordinate (C) at (3,0);
\coordinate (D) at (6,0);
\coordinate (E) at (2,0);
\draw[black!45] (-0.6,0)--(6.6,0);
\draw[blue,thick] (A)--(B);
\draw[teal!80!black,thick] (A)--(C);
\draw[black,thick] (B)--(C);
\draw[purple,very thick] (A)--(D);
\fill (A) circle (1.5pt) node[above=4pt] {$A$};
\fill (B) circle (1.5pt) node[below left=4pt] {$B$};
\fill (C) circle (1.5pt) node[below right=4pt] {$C$};
\fill (D) circle (1.5pt) node[below right=4pt] {$D$};
\draw[orange!85!black,very thick] (A)--(E);\fill (E) circle (1.5pt) node[below=4pt] {$E$};
\path (A)--(B) node[midway,above left=3pt] {$AB=4$};
\path (A)--(C) node[midway,above right=3pt] {$AC=2$};
\path (C)--(D) node[midway,above=6pt] {$CD=3$};
\end{scope}
\draw[orange!85!black,|-|] ($(B)+(0,-1.0)$)--($(E)+(0,-1.0)$) node[midway,below=3pt] {$BE=2$};
\draw[orange!85!black,|-|] ($(E)+(0,-1.0)$)--($(C)+(0,-1.0)$) node[midway,below=3pt] {$EC=1$};
\draw[purple,|-|] ($(E)+(0,-2.0)$)--($(D)+(0,-2.0)$) node[midway,below=3pt] {$DE=4$};

\end{tikzpicture}
\end{center}

\textbf{内分用和：} $BE/EC=2$，$BE+EC=3$，故 $BE=2,EC=1$。

\textbf{外分用差：} $AB>AC$，顺序为 $B,C,D$；设 $CD=y>0$，则 $BD=y+3$。
\[
 \frac{y+3}{y}=2\quad\Longrightarrow\quad y=3,\quad BD=6.
\]
于是顺序是 $B,E,C,D$，所以
\[
 \boxed{DE=EC+CD=1+3=4}.
\]
\textbf{另一算法回验：} $DE=BD-BE=6-2=4$。
三边也满足 $|4-2|<3<4+2$，题设可实现。
\end{examplebox}
